% BHU PYQ -- Manually transcribed & error-corrected
\documentclass[11pt,a4paper]{article}

\usepackage[a4paper,top=2.5cm,bottom=2.5cm,left=2.8cm,right=2.8cm,headheight=14pt]{geometry}
\usepackage{amsmath,amssymb}
\usepackage[T1]{fontenc}
\usepackage[utf8]{inputenc}
\usepackage{lmodern}
\usepackage{microtype}
\usepackage{fancyhdr}
\usepackage{enumitem}
\usepackage{hyperref}
\usepackage{lastpage}
\usepackage{parskip}

\hypersetup{colorlinks=true,urlcolor=black,linkcolor=black,
  pdftitle={CHB-501: Analytical Chemistry-I -- BHU PYQ},pdfauthor={Banaras Hindu University}}

\pagestyle{fancy}\fancyhf{}
\renewcommand{\headrulewidth}{0.4pt}\renewcommand{\footrulewidth}{0.4pt}
\fancyhead[L]{\small   \textit{Banaras Hindu University}}
\fancyhead[R]{\small   \textit{CHB-501: Analytical Chemistry-I}}
\fancyfoot[C]{\small Page  \thepage\ of \pageref{LastPage}}


\newlist{parts}{enumerate}{2}
\setlist[parts,1]{label=(\alph*),leftmargin=*,itemsep=5pt,topsep=3pt}
\setlist[parts,2]{label=(\roman*),leftmargin=*,itemsep=3pt,topsep=2pt}

\newcommand{\pts}[1]{\hfill   \textit{[#1]}}


\begin{document}


\begin{center}
  {\large\bfseries BANARAS HINDU UNIVERSITY}\\[2pt]
  {\normalsize B.Sc. (Hons.) Semester V Examination, 2013-14}\\[6pt]
  \rule{0.8\linewidth}{0.5pt}\\[6pt]
  {\large\bfseries Paper No. CHB-501: Analytical Chemistry-I}\\[4pt]
  {\normalsize Time: 3 Hours\hfill Maximum Marks: 70}
\end{center}
\rule{\linewidth}{0.4pt}
\small



\noindent   \textit{Instructions: Answer five questions in total, including Question~1 which is compulsory.}

\normalsize
\bigskip




\noindent   \textbf{Question 1.} Answer all parts of the following: \pts{5 $ \times$ 2 = 10}

\begin{parts}
  \item Differentiate between metaperiodate ($ \text{IO}_4^-$) and paraperiodate ($ \text{H}_5 \text{IO}_6$) structures.
  \item Define absolute error and relative error.
  \item What are indicators in complexometric titrations?
  \item Discuss the function of starch as an indicator in iodometry.
  \item What is the role of masking agents in analytical chemistry?
\end{parts}

\medskip\rule{\linewidth}{0.2pt}




\noindent   \textbf{Question 2.}

\begin{parts}
  \item Discuss the classification and minimization of errors in chemical analysis. \pts{6}
  \item What do you mean by significant figures? Apply computation rules to the following calculations so that significant figures are retained in the final results: \pts{8}
  
\begin{parts}
    \item $3.4 + 0.020 + 7.31 = ?$
    \item $(2.4  \times 4.02) / 0.965 = ?$
    \item $\log(4.000  \times 10^{-5}) = ?$
    \item $ \text{antilog}(12.5) = ?$
  \end{parts}
\end{parts}

\medskip\rule{\linewidth}{0.2pt}




\noindent   \textbf{Question 3.}

\begin{parts}
  \item Name the titrations which you recommend for the quantitative evaluation of organic functional groups in the following reactions, and write chemical equations: \pts{8}
  
\begin{parts}
    \item Acetylation of primary alcohols using acetic anhydride
    \item Estimation of aldehydes using hydroxylamine hydrochloride
  \end{parts}
  \item Discuss the theoretical and practical aspects of Karl Fischer reagent for the estimation of water content. \pts{6}
\end{parts}

\medskip\rule{\linewidth}{0.2pt}




\noindent   \textbf{Question 4.}

\begin{parts}
  \item Discuss the structural features of EDTA as a chelating agent in complexometric titrations. \pts{6}
  \item Explain with specific examples: \pts{8}
  
\begin{parts}
    \item EDTA displacement titration
    \item EDTA back-titration
    \item Metallochromic indicator transition mechanisms
  \end{parts}
\end{parts}

\medskip\rule{\linewidth}{0.2pt}




\noindent   \textbf{Question 5.} Write explanatory notes on any two of the following: \pts{2 $ \times$ 7 = 14}

\begin{parts}
  \item Global warming and its mitigation strategies
  \item Ozone depletion mechanism in the stratosphere
  \item Bhopal Gas Tragedy and industrial disaster management
\end{parts}

\vfill
\rule{\linewidth}{0.4pt}

\begin{center}\small
  \href{https://bhu-exam.vercel.app/}{Click here to visit our website}\\[4pt]
  \href{https://me-aryan.vercel.app/}{Click here to visit our contributor}\\[4pt]
  \href{https://quashan-me.vercel.app/}{Click here to visit our contributor}
\end{center}

\end{document}
